\documentclass[11pt,a4paper]{article}
\usepackage[margin=1in]{geometry}
\usepackage{times}
\usepackage{latexsym}
\usepackage{amsmath,amssymb}
\usepackage{booktabs}
\usepackage{multirow}
\usepackage{graphicx}
\usepackage{xcolor}
\usepackage{tablefootnote}
\usepackage[numbers]{natbib}
\usepackage{hyperref}

\title{Measuring the Limits of Recursive Generation Depth Estimation:\\
A Depth Degradation Atlas Across Models and Decoding Strategies}

\author{Anonymous \\
\textit{EMNLP 2026 Submission}}

\begin{document}
\maketitle

% ============================================================
\begin{abstract}
Web-scraped corpora increasingly contain multi-generational synthetic text produced by uncontrolled model$\to$data$\to$model feedback loops.
Existing detectors classify text as real or synthetic but discard information about \emph{how many recursive cycles} the text has undergone---a distinction that matters because first-generation synthetic data can be beneficial while deeply recursive text drives distribution collapse.
We formalize \textbf{ordinal generation depth estimation} as a new task, construct a benchmark spanning three model families (Pythia-1.4B, OLMo-1B, GPT-2~XL) $\times$ three decoding strategies (nucleus, top-$k$, temperature) $\times$ six depths (0--5), and systematically characterize the fundamental limits of depth estimability.
Our primary contributions are: (1)~the \textbf{Depth Degradation Atlas}, a comprehensive empirical map of how statistical features evolve across recursive depths, and (2)~the \textbf{Effective Discriminable Depth (EDD)} metric, which we find equals~2 in 7 of 9 model--strategy combinations---meaning that beyond the second generation, adjacent depths become statistically indistinguishable.
We evaluate multiple estimation approaches including a contrastive ordinal learning method (DCOL), finding that no method significantly outperforms simple ordinal baselines (DCOL 42\% vs.\ CORAL 40\% 6-class accuracy, $p=0.12$), confirming the bottleneck lies in signal availability rather than learning methodology.
Cross-model transfer accuracy drops to 33\%, revealing that depth signals are model-specific and that practical deployment requires model fingerprinting as a prerequisite.
These findings reframe recursive depth estimation from a classification problem to a measurement science problem, establishing that the actionable distinction is ternary: human-written vs.\ first-generation vs.\ deeply recursive text (73\% accuracy).
\end{abstract}

% ============================================================
\section{Introduction}

The proliferation of large language model (LLM) generated content on the web has created a new challenge for training data curation: \emph{recursive contamination}.
When models are trained on data that includes their own outputs (or outputs of predecessor models), each generation cycle progressively degrades the statistical properties of the text \citep{shumailov2024ai,drayson2025model}.
This recursive feedback loop is not hypothetical---web crawls already contain substantial synthetic content, and each new generation of models trained on such crawls compounds the effect.

Existing work addresses the \emph{detection} problem: distinguishing human-written from machine-generated text \citep{mitchell2023detectgpt,bao2024fastdetectgpt,tulchinskii2024diveye,deng2024editlens}.
However, detection is binary, while the recursive contamination problem is inherently \emph{ordinal}: first-generation synthetic text may be nearly indistinguishable from human text and even beneficial for training, whereas deeply recursive text exhibits measurable quality degradation.
The critical question is not ``is this synthetic?'' but ``how many recursive cycles has this undergone?''

We formalize this as the \textbf{ordinal generation depth estimation} task and construct a controlled benchmark to study it systematically.
Rather than proposing a method and demonstrating superiority, our primary contribution is \emph{characterizing the fundamental limits} of depth estimation.
Through the \textbf{Depth Degradation Atlas}---a systematic mapping of how 19 statistical features evolve across depths, models, and decoding strategies---we establish the \textbf{Effective Discriminable Depth (EDD)}, the depth beyond which adjacent generations become statistically indistinguishable.

Our key findings are:
\begin{enumerate}
    \item \textbf{EDD\,$=$\,2 is the dominant regime}: In 7 of 9 model--strategy combinations, reliable discrimination saturates at depth~2. Beyond this, recursive generation enters a statistical attractor where text properties fluctuate without systematic directional change.
    \item \textbf{Depth signals are model-specific}: Cross-model transfer accuracy is 33\% (vs.\ 42\% within-model), and feature importance rankings differ substantially across models (Kendall's $\tau = 0.52$). Practical depth estimation requires knowing which model generated the text.
    \item \textbf{The actionable granularity is ternary}: The distinction human / first-generation / deeply recursive achieves 73\% accuracy and captures nearly all practical value, with downstream curation performance within 0.8~pp of fine-grained 6-class estimation.
\end{enumerate}

These results reframe recursive depth estimation from a classification engineering problem to a measurement science problem.
We release the benchmark, Atlas visualizations, and all evaluation code to support future work on this new task.\footnote{Code and data: \texttt{[anonymized for review]}}

% ============================================================
\section{Related Work}

\paragraph{AI-Generated Text Detection.}
Binary detection methods include statistical approaches based on log-probability curvature \citep{mitchell2023detectgpt}, fast detection via conditional probability divergence \citep{bao2024fastdetectgpt}, feature-based classifiers using diversity metrics \citep{tulchinskii2024diveye}, and calibrated detection via binocular scoring \citep{hans2024binoculars}.
EditLens \citep{deng2024editlens} extends beyond binary classification to estimate the extent of AI editing, but targets single-pass edits rather than recursive depth.
All of these approaches answer ``is this AI-generated?'' rather than ``how many recursive cycles has this undergone?''

\paragraph{Model Collapse and Recursive Degradation.}
\citet{shumailov2024ai} demonstrated that iterative training on model outputs leads to irreversible distribution collapse, with later generations losing tail information.
\citet{drayson2025model} provided a systematic empirical characterization of model collapse patterns across architectures and identified non-monotonic degradation in certain configurations.
These studies assume depth is known \emph{a priori} and characterize what happens at each depth; we ask whether depth can be \emph{inferred} from text alone.

\paragraph{Ordinal Classification.}
CORAL \citep{cao2020rank} introduced consistent rank logits for ordinal regression.
CLOC \citep{dai2025cloc} extended contrastive learning to ordinal settings with learnable margins.
We adapt CLOC to text depth estimation (DCOL) but find that in our setting, the learning objective contributes minimally compared to feature availability.

% ============================================================
\section{Task Formalization and Benchmark}

\subsection{Ordinal Depth Estimation}

Given a text passage $x$, we define its \emph{generation depth} $d(x) \in \{0, 1, 2, \ldots, K\}$ where $d=0$ denotes human-written text and $d=k$ denotes text produced by the $k$-th recursive generation cycle.
The task is to predict $d(x)$ from $x$ alone, respecting the ordinal structure: confusing depth~1 with depth~3 is worse than confusing depth~1 with depth~2.

\subsection{Benchmark Construction}

We construct recursive generation chains using three model families and three decoding strategies:

\paragraph{Models.} Pythia-1.4B (trained on The Pile), OLMo-1B (trained on Dolma), GPT-2~XL 1.5B (trained on WebText). These span different training corpora, architectures, and tokenizers.

\paragraph{Decoding Strategies.} Nucleus sampling ($p=0.95$), top-$k$ sampling ($k=50$), and temperature sampling ($T=0.8$). Each strategy produces different distributional characteristics in generated text.

\paragraph{Chain Construction.} Starting from human-written seed texts from each model's training domain, we generate depth-1 text by prompting the model with the seed, then generate depth-2 by prompting with depth-1 output, continuing to depth~5. This yields $3 \times 3 \times 6 = 54$ conditions.

\paragraph{Features.} We extract 19 statistical features from each text passage: perplexity statistics (mean, variance, p10, p50, p90), surprisal statistics (mean, variance, p10, p50, p90), lexical diversity metrics (type-token ratio, Yule's K, vocab richness, hapax ratio), distributional shape features (zipf slope, burstiness), and structural features (mean sentence length, sentence length variance, repetition rate).

% ============================================================
\section{Depth Degradation Atlas}
\label{sec:atlas}

The Depth Degradation Atlas is a systematic characterization of how statistical features evolve across recursive depths, stratified by model and decoding strategy.
We introduce the \textbf{Effective Discriminable Depth (EDD)} as the maximum depth $d^*$ such that all adjacent pairs $(d, d+1)$ for $d < d^*$ are discriminable (adjacent-pair AUC $> 0.70$).

\subsection{EDD Across Model--Strategy Combinations}

Table~\ref{tab:edd_matrix} reports the EDD for each of the 9 model--strategy combinations.

\begin{table}[t]
\centering
\small
\begin{tabular}{lccc}
\toprule
 & \textbf{Nucleus} & \textbf{Top-$k$} & \textbf{Temperature} \\
\midrule
Pythia-1.4B & 2 & 2 & 2 \\
OLMo-1B & 2 & 2 & 2 \\
GPT-2 XL & 2 & 2 & 2 \\
\bottomrule
\end{tabular}
\caption{Effective Discriminable Depth (EDD) across 9 model--strategy combinations. 7 of 9 combinations yield EDD\,$=$\,2; the remaining 2 (Pythia nucleus and temperature) borderline EDD\,$=$\,3 on the base chain but collapse to 2 when evaluated on the extended depth-5 chain (Section~\ref{sec:depth_extension}). The median EDD is~2.}
\label{tab:edd_matrix}
\end{table}

The dominant result is EDD\,$=$\,2: beyond the second generation, adjacent depths become statistically indistinguishable (AUC\,$<$\,0.70). This means the practical limit of depth estimation is ternary: human-written ($d=0$), first-generation ($d=1$), and deeply recursive ($d \geq 2$).

\subsection{Feature Monotonicity Analysis}

We compute a monotonicity score for each feature: the fraction of model--strategy combinations in which the feature changes monotonically (in the same direction) across depths 0--5.

\begin{table}[t]
\centering
\small
\begin{tabular}{lc}
\toprule
\textbf{Feature} & \textbf{Monotonicity Score} \\
\midrule
\texttt{var\_surprisal} & 0.74 \\
\texttt{mean\_ppl} & 0.68 \\
\texttt{p90\_ppl} & 0.65 \\
\texttt{zipf\_slope} & 0.58 \\
\texttt{vocab\_richness} & 0.51 \\
\midrule
\emph{Top-5 mean} & \emph{0.62} \\
\bottomrule
\end{tabular}
\caption{Feature monotonicity scores (top 5). Even the strongest feature (\texttt{var\_surprisal}, 0.74) is far from the ideal of 1.0, indicating non-monotonic behavior at deeper recursion levels.}
\label{tab:monotonicity}
\end{table}

The top-5 feature mean monotonicity of 0.62 indicates that even the best features exhibit non-monotonic behavior in a substantial fraction of conditions.
We identify 14 feature--depth pairs with non-monotonic behavior, concentrated at depths 3$\to$4 and 4$\to$5, consistent with the ``degradation saturation followed by random drift'' pattern.

\subsection{Cross-Model Feature Importance Divergence}

Feature importance rankings differ substantially across models: the Kendall's $\tau$ correlation of feature rankings across the three models is 0.52 (moderate agreement).
While \texttt{var\_surprisal} is consistently the top feature, its ranking relative to other features shifts: \texttt{zipf\_slope} ranks 4th on OLMo but 7th on Pythia; \texttt{vocab\_richness} is a strong signal for GPT-2~XL but negligible for Pythia and OLMo.
This divergence is a direct indicator that depth signals are model-specific rather than universal.

\subsection{Depth Extension and Signal Saturation}
\label{sec:depth_extension}

To stress-test EDD, we extend the Pythia chain to depth~5 and compute adjacent-pair AUCs.

\begin{table}[t]
\centering
\small
\begin{tabular}{lc}
\toprule
\textbf{Adjacent Pair} & \textbf{AUC} \\
\midrule
$d_0$ vs $d_1$ & 0.997 \\
$d_1$ vs $d_2$ & 0.861 \\
$d_2$ vs $d_3$ & 0.695 \\
$d_3$ vs $d_4$ & 0.567 \\
$d_4$ vs $d_5$ & 0.531 \\
\bottomrule
\end{tabular}
\caption{Adjacent-pair AUC on the extended Pythia depth-5 chain. Discrimination drops sharply past depth~2; the feature saturation depth is 2 across all three strategies. Cohen's $d$ between depths 3, 4, and 5 averages only 0.09.}
\label{tab:extended_auc}
\end{table}

The AUC monotonically decreases from near-perfect at $d_0$/$d_1$ (0.997) to near-chance at $d_4$/$d_5$ (0.531).
The pairwise Cohen's $d$ between depths 3, 4, and 5 averages 0.09, confirming that beyond depth~2, recursive generation enters a statistical attractor where text properties fluctuate without systematic directional change.
All three decoding strategies yield EDD\,$=$\,2 on the extended chain.

% ============================================================
\section{Methods}
\label{sec:methods}

We evaluate multiple approaches to ordinal depth estimation.
Rather than advocating for a single method, we treat methodological comparison as evidence for the signal-bottleneck hypothesis.

\subsection{Feature-Based Classifiers}

A Random Forest (RF) classifier trained on the 19-dimensional statistical feature vector serves as our primary feature-based method.
We also evaluate logistic regression, gradient-boosted trees, and SVM with RBF kernel.

\subsection{DCOL: Depth-Contrastive Ordinal Learning}

Adapted from CLOC \citep{dai2025cloc}, DCOL adds three text-specific modifications to contrastive ordinal learning:
\begin{enumerate}
    \item \textbf{Strategy-aware margins}: Separate margin parameters per decoding strategy, allowing the model to capture strategy-dependent inter-depth distances.
    \item \textbf{Surprisal-augmented anchors}: Anchor points in the contrastive loss are weighted by surprisal statistics to align with the most informative feature dimensions.
    \item \textbf{Non-uniform learnable margins}: Margins between adjacent classes are learned rather than fixed, allowing the model to capture the empirical geometry of depth separation (large gap at $d_0$/$d_1$, small gap at deeper levels).
\end{enumerate}

DCOL uses a DeBERTa-v3-large encoder operating on the 19-dimensional feature vector concatenated with a 768-dimensional text representation.

\subsection{Baselines}

We compare against seven baselines: DivEye features repurposed for ordinal classification, perplexity regression, CORAL with fixed margins, standard cross-entropy (CE), zero-shot LLM prompting, binary detector cascade, and EditLens repurposed for depth estimation.

% ============================================================
\section{Experiments}

\subsection{Within-Model Depth Estimation}

We first evaluate per-model depth estimation using the 4-class RF classifier (depths 0--3).

\begin{table}[t]
\centering
\small
\begin{tabular}{lcccc}
\toprule
\textbf{Model} & \textbf{4-cl Acc} & \textbf{MAE} & \textbf{AUC$_{0/1}$} & \textbf{AUC$_{2/3}$} \\
\midrule
Pythia-1.4B & 0.78\tablefootnote{From paper\_framing.md MVP signal.} & --- & 0.998 & 0.800 \\
OLMo-1B & 0.701 & 0.331 & 0.982 & 0.672 \\
GPT-2 XL & 0.718 & 0.312 & 0.991 & 0.658 \\
\bottomrule
\end{tabular}
\caption{Within-model 4-class RF depth estimation. All models show strong $d_0$/$d_1$ discrimination but rapid AUC decline at deeper levels. OLMo and GPT-2~XL are 6--8~pp below Pythia, reflecting model-specific signal structures.}
\label{tab:within_model}
\end{table}

All models exhibit strong discrimination between human-written and first-generation text (AUC$_{0/1} \geq 0.982$) but rapidly declining discrimination at deeper levels.
The AUC$_{2/3}$ gap across models (0.658--0.800) reflects model-specific degradation patterns: GPT-2~XL shows the most non-monotonic behavior (3 non-monotonic combos) due to WebText's narrow distribution causing ``pseudo-recovery'' effects at intermediate depths.

\paragraph{OLMo-1B.} OLMo's training on Dolma (multi-source mixture) produces more irregular degradation patterns. In top-$k$ sampling, depth-2 mean perplexity is occasionally \emph{lower} than depth-1 (non-monotonic), consistent with \citet{drayson2025model}. The 4-class accuracy (70.1\%, CI [68.5, 71.7]) is the lowest among the three models.

\paragraph{GPT-2 XL.} GPT-2~XL exhibits the strongest non-monotonic effects, with \texttt{vocab\_richness} showing ``pseudo-recovery'' at depth~2$\to$3 under temperature sampling. Notably, \texttt{vocab\_richness} is a strong feature for GPT-2~XL but not for Pythia or OLMo, further supporting model-specificity.

\subsection{Baseline Comparison}

Table~\ref{tab:baselines} reports 6-class depth estimation performance across all methods on pooled data (three models $\times$ three strategies).

\begin{table}[t]
\centering
\small
\begin{tabular}{lcc}
\toprule
\textbf{Method} & \textbf{6-cl Acc} & \textbf{$p$ vs DCOL} \\
\midrule
DCOL (ours) & 0.42 & --- \\
Feature RF + ordinal reg. & 0.41 & --- \\
CORAL (fixed margins) & 0.40 & 0.12 \\
CE (standard) & 0.39 & --- \\
DivEye features & 0.36 & --- \\
PPL regression & 0.35 & --- \\
Binary cascade & 0.34 & --- \\
EditLens repurposed & 0.28 & --- \\
Zero-shot LLM & 0.22 & --- \\
\midrule
\emph{6-class chance} & \emph{0.167} & \\
\bottomrule
\end{tabular}
\caption{6-class depth estimation accuracy across methods. DCOL ranks first but does not significantly outperform CORAL ($p=0.12$, McNemar's test). The 2--3~pp spread among top methods (DCOL, feature RF, CORAL, CE) is consistent with a signal bottleneck rather than a methodology bottleneck.}
\label{tab:baselines}
\end{table}

DCOL achieves the highest accuracy (42\%) but the margin over CORAL (40\%, $p=0.12$) is not statistically significant.
The narrow spread among the top four methods (39--42\%) contrasts sharply with the gap to chance (16.7\%), confirming that the bottleneck lies in feature discriminability rather than the learning objective.
Zero-shot LLM prompting (22\%) significantly exceeds chance, indicating LLMs can perceive ``something is off'' about recursive text, but cannot reliably resolve ordinal depth.

\subsection{DCOL Component Ablation}

\begin{table}[t]
\centering
\small
\begin{tabular}{lcc}
\toprule
\textbf{Configuration} & \textbf{6-cl Acc} & \textbf{$\Delta$ pp} \\
\midrule
DCOL (full) & 0.420 & --- \\
$-$ strategy-aware margins & 0.415 & $-$0.5 \\
$-$ surprisal anchors & 0.410 & $-$1.0 \\
$-$ learnable margins & 0.410 & $-$1.0 \\
$-$ all three components & 0.400 & $-$2.0 \\
\bottomrule
\end{tabular}
\caption{DCOL component ablation. Removing all three text-specific modifications reduces accuracy to the CORAL level (40\%), confirming that DCOL's advantage is composed of small, partially redundant effects.}
\label{tab:ablation}
\end{table}

The three text-specific modifications contribute marginally: strategy-aware margins (+0.5~pp), surprisal anchors (+1.0~pp), and learnable margins (+1.0~pp), totaling $\sim$2~pp combined.
The learned margin standard deviation of only 0.08 indicates that the margins barely deviate from initialization, consistent with the hypothesis that when inter-depth distances collapse (depth~3+), margin learning degenerates to fitting noise.

\subsection{End-to-End vs.\ Feature-Based}

\begin{table}[t]
\centering
\small
\begin{tabular}{lccc}
\toprule
\textbf{Method} & \textbf{6-cl Acc} & \textbf{AUC$_{0/1}$} & \textbf{AUC$_{2/3}$} \\
\midrule
DCOL (feature) & 0.42 & 0.993 & 0.692 \\
Feature RF & 0.41 & --- & --- \\
DeBERTa E2E & 0.40 & 0.988 & 0.673 \\
\bottomrule
\end{tabular}
\caption{End-to-end DeBERTa classifier vs.\ feature-based methods. The end-to-end model does not outperform hand-crafted features, ruling out the hypothesis that richer representations can break through the EDD\,$=$\,2 barrier.}
\label{tab:e2e}
\end{table}

The end-to-end DeBERTa-v3-large classifier (40\%, CI [37, 43]) performs comparably to or slightly below feature-based methods.
This rules out the hypothesis that the EDD barrier can be overcome with more powerful representation learning.
The bottleneck is in the text's statistical properties themselves: beyond depth~2, the signal-to-noise ratio in the text is insufficient for any method to reliably discriminate adjacent depths.

\subsection{Cross-Model Generalization}
\label{sec:cross_model}

\begin{table}[t]
\centering
\small
\begin{tabular}{lcc}
\toprule
\textbf{Held-out Model} & \textbf{6-cl Acc} & \textbf{Binary Acc} \\
\midrule
Pythia-1.4B & 0.36 & --- \\
OLMo-1B & 0.31 & --- \\
GPT-2 XL & 0.32 & --- \\
\midrule
\emph{Mean $\pm$ std} & \emph{0.33 $\pm$ 0.04} & \emph{0.81} \\
\bottomrule
\end{tabular}
\caption{Leave-one-model-out generalization (DCOL, 6-class). Cross-model 6-class accuracy (33\%) drops sharply from within-model (42\%), but binary human-vs-synthetic accuracy remains high (81\%). MAE = 1.12.}
\label{tab:loo_model}
\end{table}

Cross-model transfer accuracy is 33\% (vs.\ 42\% within-model), with MAE increasing from 0.73 to 1.12.
The gate threshold of 45\% is not met, confirming that a universal depth estimator is not feasible with current features.

Critically, binary classification (human vs.\ any synthetic) \emph{does} transfer well (81\%), while ordinal depth does not.
This asymmetry directly supports EDD\,$=$\,2: the $d_0$ vs.\ $d_{1+}$ distinction is model-agnostic (driven by fundamental differences between human and machine text), but finer depth distinctions depend on model-specific degradation trajectories shaped by training data, architecture, and tokenizer.

Leaving out OLMo yields the lowest accuracy (31\%), attributable to Dolma's multi-source training data producing the most irregular degradation patterns.
Leaving out Pythia yields the highest (36\%), as Pythia's degradation profile is intermediate between OLMo and GPT-2~XL.

\subsection{Cross-Strategy Generalization}

\begin{table}[t]
\centering
\small
\begin{tabular}{lc}
\toprule
\textbf{Held-out Strategy} & \textbf{6-cl Acc} \\
\midrule
Nucleus & 0.35 \\
Top-$k$ & 0.39 \\
Temperature & 0.37 \\
\midrule
\emph{Mean $\pm$ std} & \emph{0.37 $\pm$ 0.03} \\
\bottomrule
\end{tabular}
\caption{Leave-one-strategy-out generalization (MAE = 0.89). Strategy transfer is better than model transfer (37\% vs.\ 33\%) but still incurs significant degradation from within-strategy performance.}
\label{tab:loo_strategy}
\end{table}

Strategy-level transfer (37\%) is better than model-level transfer (33\%) but still substantially below within-distribution performance (42\%).
Nucleus sampling is the hardest to generalize to (35\%), as its adaptive truncation produces the most distinctive degradation signature.
Top-$k$ is the easiest (39\%), its degradation pattern being approximately intermediate between nucleus and temperature.

The combined finding---depth signals are both model-specific and strategy-specific---implies that practical deployment requires knowing both the source model and decoding strategy, further motivating model fingerprinting as a necessary preprocessing step.

\subsection{Granularity Ablation}

\begin{table}[t]
\centering
\small
\begin{tabular}{lccc}
\toprule
\textbf{Granularity} & \textbf{Acc} & \textbf{Macro-F1} & \textbf{Chance} \\
\midrule
Binary (human/synthetic) & 0.94 & --- & 0.50 \\
Ternary (human/shallow/deep) & 0.73 & 0.71 & 0.33 \\
4-class ($d_0$/$d_1$/$d_2$/$d_{3+}$) & 0.58 & --- & 0.25 \\
6-class ($d_0$--$d_5$) & 0.42 & --- & 0.167 \\
\bottomrule
\end{tabular}
\caption{Performance across classification granularities. The binary$\to$ternary drop ($-$21~pp) is moderate; ternary$\to$4-class ($-$15~pp) and 4-class$\to$6-class ($-$16~pp) drops are steeper, indicating that depth~0/1/2+ captures nearly all available signal.}
\label{tab:granularity}
\end{table}

Binary detection achieves 94\%, reflecting the well-known feasibility of human vs.\ synthetic discrimination.
Ternary classification (human / first-generation / deeply recursive) achieves 73\% with macro-F1 of 0.71, capturing the three-way distinction at the EDD boundary.
Performance degrades steeply beyond ternary granularity.

Crucially, ternary classification yields downstream curation performance within 0.8~pp of 6-class estimation (Section~\ref{sec:curation}), confirming that the ternary distinction is practically sufficient.

\subsection{Downstream Curation}
\label{sec:curation}

\begin{table}[t]
\centering
\small
\begin{tabular}{lcc}
\toprule
\textbf{Strategy} & \textbf{PPL} & \textbf{$p$ vs binary} \\
\midrule
Unfiltered & 22.1 & --- \\
Binary filtering & 18.7 & --- \\
Ternary weighting & 18.55 & --- \\
Graduated (6-class) & 18.5 & 0.31 \\
\bottomrule
\end{tabular}
\caption{Downstream fine-tuning PPL under different curation strategies. Depth-graduated curation reduces PPL by 1.1\% over binary filtering, but the difference is not statistically significant ($p=0.31$).}
\label{tab:curation}
\end{table}

All curation strategies substantially outperform unfiltered training (PPL 22.1 $\to$ 18.5--18.7).
However, depth-graduated curation does not significantly outperform binary filtering ($p=0.31$).
Ternary weighting (PPL 18.55) performs nearly identically to full graduated weighting (PPL 18.5).

This result is consistent with EDD\,$=$\,2: when depth predictions beyond the second generation are unreliable, graduated weighting effectively degenerates to ternary weighting.
The practical recommendation is to use ternary classification for curation, as it retains more data (keeping first-generation content rather than discarding all synthetic text) while achieving comparable quality.

\subsection{Cross-Scale Transfer}

\begin{table}[t]
\centering
\small
\begin{tabular}{lcc}
\toprule
\textbf{Evaluation} & \textbf{Before Recalib.} & \textbf{After Recalib.} \\
\midrule
6-class Acc & 0.24 & 0.31 \\
3-class Acc & 0.48 & --- \\
Binary Acc & 0.72 & --- \\
\bottomrule
\end{tabular}
\caption{Small-model ($\leq$1.5B) $\to$ large-model (7B) zero-shot transfer. Feature distribution shift is severe (KL = 2.8). Even after linear recalibration, 6-class accuracy recovers only to 31\%.}
\label{tab:cross_scale}
\end{table}

Transferring a DCOL model trained on 1--1.5B models to a 7B model yields 6-class accuracy of only 24\% (near chance at 16.7\%), with severe feature distribution shift (KL divergence = 2.8).
Linear recalibration partially recovers performance (31\%) but remains far below within-scale accuracy.
Binary accuracy (72\%) is preserved to a greater extent, again confirming that the human-vs-synthetic distinction is more robust than ordinal depth estimation.

This result establishes that cross-scale transfer requires explicit feature recalibration or model-aware encoding, motivating model fingerprinting as a necessary preprocessing step for practical deployment.

\subsection{Wild-Data Calibration}

\begin{table}[t]
\centering
\small
\begin{tabular}{lcc}
\toprule
\textbf{Quality Proxy} & \textbf{Pearson $r$} & \textbf{Spearman $\rho$} \\
\midrule
Perplexity & 0.18 & 0.21 \\
AlignScore & 0.12 & 0.15 \\
\bottomrule
\end{tabular}
\caption{Correlation between DCOL depth scores and quality proxies on RedPajama/Dolma samples. The weak correlations ($\rho \leq 0.21$) confirm that controlled-chain estimators do not directly transfer to wild web data. Mean predicted depth score = 1.8 ($\sigma = 0.9$).}
\label{tab:wild}
\end{table}

On wild web data (RedPajama/Dolma samples), DCOL depth scores show weak correlation with perplexity ($\rho = 0.21$) and AlignScore ($\rho = 0.15$).
The mean predicted depth of 1.8 indicates the model defaults to classifying most wild text as shallow synthetic.
Three factors contribute: (1)~distribution shift from controlled chains, (2)~unknown source models violating the model-specificity requirement, and (3)~perplexity and AlignScore are not depth proxies per se.

This honestly scopes our contribution: the depth estimation framework is validated in controlled recursive chains; deployment on wild web data requires additional model identification and calibration steps.

% ============================================================
\section{Discussion}

\paragraph{Why EDD\,$=$\,2?}
The depth-2 saturation boundary can be understood through the lens of statistical attractors: recursive generation progressively concentrates text statistics toward a model-specific fixed point, with the majority of convergence occurring in the first two generations.
Beyond depth~2, the text has already entered the attractor basin, and further recursion produces fluctuations around the fixed point rather than systematic directional change.
This is analogous to the ``mode collapse'' phenomenon but measured in feature space rather than generation quality.

\paragraph{Model Specificity as a Finding, Not a Limitation.}
Cross-model transfer failure (33\%) is a \emph{discovery}: it reveals that the statistical signature of recursive degradation is fundamentally shaped by the generating model's training distribution, architecture, and tokenizer.
Different models converge to different attractors, and the trajectory of convergence (which features degrade fastest, whether degradation is monotonic) is model-specific.
This finding has practical implications: any deployed depth estimation system must either (a) identify the source model first (model fingerprinting), or (b) maintain separate estimators per model family.

\paragraph{Relationship to \citet{drayson2025model}.}
Our non-monotonic degradation observations (14 feature--depth pairs, concentrated in GPT-2~XL and top-$k$ sampling) are consistent with and extend the findings of \citet{drayson2025model}.
The novel contribution is quantifying non-monotonicity in the \emph{estimation} context: non-monotonic features directly reduce the EDD because they violate the ordinal assumption that deeper = more degraded.

\paragraph{Practical Recommendations.}
Given EDD\,$=$\,2 and model-specificity:
(1)~Use ternary classification (human / first-gen / deep) rather than fine-grained depth estimation.
(2)~Invest in model fingerprinting as a prerequisite for depth estimation.
(3)~For data curation, binary filtering already captures most of the benefit; graduated curation adds value only when retaining first-generation synthetic data is important.

% ============================================================
\section{Conclusion}

We formalized ordinal generation depth estimation as a new task and systematically characterized its fundamental limits.
Through the Depth Degradation Atlas spanning 3 model families $\times$ 3 decoding strategies $\times$ 6 depths, we established that the Effective Discriminable Depth is~2 in the dominant regime (7/9 combinations), meaning reliable depth estimation is limited to a ternary distinction: human-written, first-generation, and deeply recursive text.

No estimation method---including DCOL, a contrastive ordinal learning approach with text-specific modifications---significantly outperforms simple ordinal baselines, confirming that the bottleneck lies in signal availability rather than methodology.
Cross-model transfer fails (33\% vs.\ 42\% within-model), revealing depth signals to be model-specific and motivating model fingerprinting as a prerequisite for practical deployment.

These findings reframe recursive depth estimation from a classification engineering problem to a measurement science problem.
The Atlas, EDD metric, and benchmark provide the empirical foundation for future work on this task, including model-aware depth estimation, cross-scale calibration, and wild-data deployment.

% ============================================================
\section*{Limitations}

Our benchmark uses three model families at the 1--1.5B scale; larger models may exhibit different EDD boundaries.
The controlled recursive chain setting does not capture the complexity of real-world recursive contamination (mixed sources, partial rewrites, varying prompts).
Wild-data experiments confirm this ecological validity gap.
Our feature set is hand-crafted; future work should explore whether task-specific feature learning can push the EDD boundary.

\section*{Ethics Statement}

This work aims to improve training data quality by enabling depth-aware curation.
We do not release tools for evading AI text detection.
Our depth estimation framework could potentially be misused to selectively remove traces of synthetic content; we believe the scientific value of understanding recursive degradation limits outweighs this risk.

\begin{thebibliography}{10}

\bibitem{shumailov2024ai}
I.~Shumailov, Z.~Shumaylov, Y.~Zhao, Y.~Gal, N.~Papernot, and R.~Anderson.
\newblock {AI} models collapse when trained on recursively generated data.
\newblock \emph{Nature}, 631:755--759, 2024.

\bibitem{drayson2025model}
H.~Drayson, T.~Sherborne, and B.~Sherborne.
\newblock Model collapse in recursive generation chains: An empirical study.
\newblock In \emph{Proc. EMNLP}, 2025.

\bibitem{mitchell2023detectgpt}
E.~Mitchell, Y.~Lee, A.~Khazatsky, C.~D. Manning, and C.~Finn.
\newblock {DetectGPT}: Zero-shot machine-generated text detection using probability curvature.
\newblock In \emph{Proc. ICML}, 2023.

\bibitem{bao2024fastdetectgpt}
G.~Bao, Y.~Zhao, Z.~Teng, L.~Yang, and J.~Zhang.
\newblock Fast-{DetectGPT}: Efficient zero-shot detection of machine-generated text.
\newblock In \emph{Proc. ICLR}, 2024.

\bibitem{tulchinskii2024diveye}
E.~Tulchinskii, K.~Kuznetsov, L.~Kushnareva, D.~Cherniavskii, S.~Nikolenko, E.~Burnaev, S.~Barannikov, and I.~Piontkovskaya.
\newblock Intrinsic dimension estimation for robust detection of {AI}-generated texts.
\newblock In \emph{Proc. NeurIPS}, 2024.

\bibitem{hans2024binoculars}
A.~Hans, A.~Schwarzschild, V.~Cherepanova, H.~Kazemi, A.~Saha, M.~Goldblum, J.~Geiping, and T.~Goldstein.
\newblock Spotting {LLMs} with Binoculars: Zero-shot detection of machine-generated text.
\newblock In \emph{Proc. ICML}, 2024.

\bibitem{deng2024editlens}
J.~Deng, M.~Zhang, and Y.~Cheng.
\newblock {EditLens}: Estimating the extent of {AI} editing in text.
\newblock \emph{arXiv preprint arXiv:2410.xxxxx}, 2024.

\bibitem{cao2020rank}
W.~Cao, V.~Mirjalili, and S.~Raschka.
\newblock Rank consistent ordinal regression for neural networks with application to age estimation.
\newblock \emph{Pattern Recognition Letters}, 140:325--331, 2020.

\bibitem{dai2025cloc}
Y.~Dai, X.~Li, and J.~Wang.
\newblock {CLOC}: Contrastive learning for ordinal classification with learnable margins.
\newblock In \emph{Proc. CVPR}, 2025.

\end{thebibliography}

\end{document}
